\documentclass[11pt]{article}
\usepackage[T1]{fontenc}
\usepackage{lmodern,amsmath,amssymb,booktabs,geometry,hyperref}
\geometry{margin=1in}
\pdfmapfile{+lm.map}
\pdfmapfile{+cm.map}
\pdfmapfile{+symbols.map}
\pdfmapfile{+latxfont.map}
\pdfmapfile{+euler.map}
\setlength{\emergencystretch}{3em}
\hypersetup{colorlinks=true,linkcolor=blue,urlcolor=blue,citecolor=blue}
\newcommand{\F}{\mathbb F_2}
\newcommand{\reg}{\operatorname{reg}}
\newcommand{\im}{\operatorname{im}}
\title{Two elementary counterexamples to\\TLMC conjecture 00000003322}
\author{AI-assisted solution prepared for TLMC}
\date{21 October 2026}
\begin{document}
\maketitle
\begin{abstract}
The asserted equivalence between regularity equaling the matching number and being a bipartite forest is false. Under the literal edge-ideal convention, the single edge $K_2$ has regularity $2$ and matching number $1$, despite being a bipartite forest. Under the alternative quotient-ring convention, the triangle $K_3$ has regularity $1$ and matching number $1$, despite not being bipartite. Both computations follow from elementary minimal graded free resolutions over $\F$, constructed and verified in a dependency-free Lean 4.31.0 project.
\end{abstract}

\section{Exact target and the normalization issue}
The numerical equivalence in the source is:
\begin{quote}
``regularity equals the matching number if and only if the graph is a bipartite forest''.
\end{quote}
The Chinese statement asserts the same equivalence. The full bilingual text is preserved in \texttt{source.md}, with upstream Git blob
\texttt{c13f3ca1b54b390533587ff148a83be0d1800bf8}. Its reference to a generalization involving Cohen--Macaulay trees does not define that generalization. We make no claim about it; our target is the explicit equivalence above.

For a finite simple graph $G$ on variables $x_1,\ldots,x_n$, its edge ideal is
\[
 I(G)=(x_ix_j:\{i,j\}\in E(G))\subseteq k[x_1,\ldots,x_n].
\]
Let $\nu(G)$ be the maximum size of a collection of vertex-disjoint edges. We use standard total-degree grading and the minimal-resolution characterization
\[
 \reg M=\max\{d-i:\beta_{i,d}(M)\ne0\}.
\]
These definitions are recalled in Sections 2.1--2.3 of \cite{survey}. Some graph literature writes graph regularity for $\reg I(G)$, while other work uses $\reg(k[V]/I(G))$; the introduction of \cite{quotient} explicitly adopts that convention. We do not silently choose between those readings: the two counterexamples below separately refute both. Both graphs are connected, so no disconnected-graph convention is involved. The formal witnesses use $k=\F$.

\section{Literal ideal regularity: the single edge}
Let $G=K_2$, with vertices $x,y$. It is bipartite, with its two singleton color classes. It is a forest because a simple cycle needs at least three distinct vertices. Its sole edge is a matching of size one, and no two edges can be vertex-disjoint. Thus $\nu(G)=1$.

In $R=\F[x,y]$, the actual edge ideal is $I=(xy)$. Multiplication by $xy$ gives a graded $R$-linear isomorphism
\[
 R(-2)\longrightarrow I,\qquad a\longmapsto xya.
\]
Surjectivity is the defining principal-ideal property. Injectivity follows directly by shifting coefficients: the coefficient of $x^{i+1}y^{j+1}$ in $xya$ is the coefficient of $x^iy^j$ in $a$. Hence the complete minimal resolution of $I$ has just $R(-2)$ in homological degree zero. Reducing modulo $(x,y)$ gives one nonzero homogeneous residue space, in degree $2$. Consequently
\[
 \reg I(K_2)=2\ne1=\nu(K_2),
\]
although $K_2$ is a bipartite forest. This disproves the right-to-left implication under the literal ideal convention.

\section{Quotient regularity: the triangle}
Let $G=K_3$, with vertices $x,y,z$. Again $\nu(G)=1$: any two edges among three vertices share a vertex. It is not bipartite, because assigning two colors to three pairwise adjacent vertices forces a monochromatic edge.

Put $S=\F[x,y,z]$ and $J=(xy,xz,yz)$. Consider
\[
 0\longrightarrow S(-3)^2\xrightarrow{D_2}S(-2)^3
 \xrightarrow{D_1}S\xrightarrow{\pi}S/J\longrightarrow0,
\]
where
\[
 D_1(a,b,c)=xya+xzb+yzc,
 \qquad
 D_2(u,v)=\bigl(z(u+v),\,yu,\,xv\bigr),
\]
and $\pi$ is the actual quotient projection. Thus the matrix of $D_2$ has columns $(z,y,0)^{\mathsf T}$ and $(z,0,x)^{\mathsf T}$. All signs are plus because the coefficient field has characteristic two.

\paragraph{Exactness, for arbitrary polynomials.}
The composition $D_1D_2$ is zero. Suppose $D_1(a,b,c)=0$. On the coefficient plane where the $x$-exponent is zero, the equality forces all coefficients of $c$ on that plane to vanish. Therefore $c=xc'$ for the polynomial $c'$ obtained by dropping one $x$ from each nonzero term. Cancelling multiplication by $x$ gives
\[
 ya+zb+yzc'=0.
\]
Examining the plane where the $y$-exponent is zero now gives $b=yb'$. Cancelling multiplication by $y$ then yields
\[
 a=z(b'+c').
\]
Thus $(a,b,c)=D_2(b',c')$, proving $\ker D_1=\im D_2$. If $D_2(u,v)=0$, its second and third coordinates give $yu=0$ and $xv=0$, so coefficient shifting gives $u=v=0$. Hence $D_2$ is injective. Finally, the image of $D_1$ is exactly the ideal generated by the three edge monomials, which is exactly $\ker\pi$, and $\pi$ is surjective.

\paragraph{Grading, minimality, and residue homology.}
Each generator of $J$ has degree two; each entry of $D_2$ has degree one. The displayed shifts therefore make both maps homogeneous of degree zero. Every matrix entry lies in the irrelevant maximal ideal $\mathfrak m=(x,y,z)$, so the resolution is minimal.

The constant-coefficient homomorphism $S\to\F$ has kernel exactly $\mathfrak m$. Applied coordinatewise, it identifies the reduced free modules with $\F$, $\F^3$, and $\F^2$, in shifts $0,2,3$, respectively. Both induced differentials vanish. Therefore the only nonzero graded homology positions of the residue complex are
\[
 (i,d)=(0,0),\quad(1,2),\quad(2,3).
\]
This gives
\[
 \reg(S/J)=\max\{0,2-1,3-2\}=1=\nu(K_3).
\]
Since $K_3$ is not bipartite, it is not a bipartite forest. This disproves the left-to-right implication under the quotient convention.

\section{What the Lean project verifies}
The project uses only Lean's bundled libraries: the mathematical modules import \texttt{Std}, while the origin-based audit also imports \texttt{Lean}. No external mathematics library or downloaded dependency is needed.
\begin{itemize}
\item \texttt{PolynomialCore}, \texttt{RawPolynomial}, \texttt{PolynomialRing}, and \texttt{TwoExtras} construct $\F[x,y]$. The \texttt{Three*} modules construct $\F[x,y,z]$. Polynomials are genuine finite-support coefficient functions, with no fixed degree bound. Multiplication is independent of the chosen finite monomial representation; all commutative-ring and coefficient-field laws are proved.
\item \texttt{Graphs} defines simple graphs, bipartitions, simple cycles, and vertex-disjoint matchings. It verifies the required properties of $K_2,K_3$, including universal upper bounds on matching size.
\item \texttt{Principal} verifies the actual $K_2$ edge ideal, the graded free rank-one resolution, the residue-field identification, and the regularity value $2$.
\item \texttt{TriangleCore} proves the unrestricted coefficient-level kernel theorem and injectivity. \texttt{TriangleModule} proves the least-ideal characterization of the $K_3$ edge ideal, constructs the actual quotient module, verifies its operations and module laws, and proves projection surjectivity and its exact kernel.
\item \texttt{TriangleGraded} proves the degree shifts, minimality, and the equality between coordinate reduction kernels and $\mathfrak m$ times the free modules. \texttt{TriangleBetti} defines the induced residue differentials, proves that they descend from the polynomial maps, and identifies nonzero homogeneous cycles modulo boundaries with the asserted support. Regularity is then computed from that support.
\item \texttt{TLMC3322.complete\_counterexamples} combines the actual resolution certificates, graph facts, both regularity values, and both false equivalences. The two convention-specific negations are separate named theorems.
\end{itemize}

Lean indexes internal degree by natural numbers. All polynomial exponents and displayed free-module shifts are nonnegative, so negative internal graded pieces vanish. The subtraction $d-i$ and regularity values are integers; no general library of arbitrarily integer-graded modules is asserted.

The proof is not a check of a proposed list of degree shifts: the maps, their exactness, the actual quotient action, grading, residue reductions, and cycles-versus-boundaries identification are all part of the formal development. The underlying polynomial claims quantify over every finitely supported polynomial, not just a search window. The project uses no admitted theorem, extra axiom, native evaluator proof, or kernel bypass. The audit allows only the usual Lean logical foundations \texttt{propext}, \texttt{Classical.choice}, and \texttt{Quot.sound}; classical choice is used to select a finite representation for multiplication.

\section{Reproduction and evidence}
With Lean 4.31.0 (including \texttt{lake} and \texttt{leanchecker}), Python 3, and a working pdfLaTeX installation already available, run \texttt{./VERIFY.sh}. It creates a fresh source-only temporary project, compiles every module with warnings as errors, prints all 323 top-level theorem types and axioms, audits all 554 theorem declarations (including generated and local theorems) by originating module, and independently kernel-replays every mathematical module. It then rebuilds this PDF without shell escape. Set \texttt{BUILD\_PDF=0} to run just the formal checks; \texttt{LEAN\_BIN\_DIR} may point to an existing Lean installation. No installation or network access occurs.

The auxiliary Python script exhaustively checks the two finite graphs and performs a clearly labeled, independent low-degree sanity check of the triangle coefficient formulas. These checks do not replace the unrestricted Lean proof. Eligibility evidence includes the live bilingual source, empty padded and unpadded PR searches, absent solution-folder response, and the false/false metadata row, together with hashes of the full PR and solution-tree snapshots used before work. Reproduction logs and the complete theorem manifest are included.

\begin{thebibliography}{9}
\bibitem{source} The Last Math Competition, conjecture 00000003322.
\href{https://github.com/The-Last-Math-Competition/The-Last-Math-Competition/blob/b355e42906e6eaf52bc7e9d005e0ef02d8aa7400/conjectures/00000003322.md}{Source pinned to commit b355e42906e6}.
\bibitem{survey} A. Banerjee, S. Kara, H. T. H\`a, \emph{Regularity of Edge Ideals and Their Powers}, Sections 2.1--2.3, arXiv:1712.00887.
\url{https://arxiv.org/abs/1712.00887}.
\bibitem{quotient} T. Biyikoglu and Y. Civan, \emph{Prime graphs, matchings and the Castelnuovo--Mumford regularity}, Journal of Commutative Algebra 11 (2019), 1--27, Introduction.
\url{https://doi.org/10.1216/JCA-2019-11-1-1}.
\end{thebibliography}
\end{document}
